% 日本語版レポート。lualatex で組む。
%   lualatex -interaction=nonstopmode hm_safety_report_ja.tex
% 英語版 hm_safety_report.tex と同じ内容・同じ数式組版。
\documentclass[11pt,a4paper]{ltjsarticle}
\usepackage{amsmath,amssymb}
\usepackage{amsthm}
\usepackage{mathpartir}
\usepackage{listings}
\usepackage{xcolor}
\usepackage{geometry}
\geometry{margin=25mm}

\newtheorem{theorem}{定理}
\newtheorem{lemma}[theorem]{補題}
\newtheorem{corollary}[theorem]{系}

\lstdefinelanguage{Coq}{
  morekeywords={
    Inductive,Definition,Fixpoint,Lemma,Theorem,Proof,Qed,forall,exists,
    match,with,end,fun,Prop,Type,Set,apply,eapply,constructor,intros,
    inversion,subst,destruct,reflexivity,exact
  },
  sensitive=true,
  morecomment=[n]{(*}{*)}
}
\lstset{
  language=Coq,
  basicstyle=\ttfamily\small,
  keywordstyle=\color{blue!60!black},
  commentstyle=\color{green!35!black},
  breaklines=true,
  frame=single,
  columns=fullflexible
}

\newcommand{\tint}{\mathtt{int}}
\newcommand{\tbool}{\mathtt{bool}}
\newcommand{\ftv}{\mathrm{ftv}}
\newcommand{\gen}{\mathrm{gen}}
\newcommand{\genmax}{\mathrm{gen}_{\max}}
\newcommand{\inst}{\preceq}
\newcommand{\subst}[3]{#1[#2 := #3]}
\newcommand{\step}{\longrightarrow}

\title{Hindley--Milner 型推論の健全性と型安全性\\\large Coq による形式化と機械検証}
\author{}
\date{}

\begin{document}
\maketitle

\section{目的}

本稿は、小さな Hindley--Milner 言語について Coq で機械検証した二つの結果を
述べる。Coq の記述と対照できるよう、文法と型付け規則を通常の数学記法で明示し、
証明の進み方も具体的に記す。

\begin{enumerate}
  \item \textbf{型推論の健全性}: 推論の判断が型を与えるなら、宣言的な判断も
  同じ型を与える。\texttt{HMSoundness.v} である。
  \item \textbf{型安全性}: 閉じた型付き項は値であるか一歩実行でき、実行しても
  型が保存される。\texttt{HMTypeSafety.v} である。
\end{enumerate}

第\ref{sec:sound-syntax}節から第\ref{sec:sound-proof}節までが前者の文法・規則・
証明、第\ref{sec:safe-syntax}節から第\ref{sec:safe-proof}節までが後者である。
覆っていない範囲は第\ref{sec:limits}節に述べる。

\section{今回の改訂で直したこと}

以前の版は Hindley--Milner を名乗りながら、その中身を持っていなかった。

\paragraph{具体化と一般化が退化していた。}
\texttt{HMSoundness.v} は \texttt{inst (Forall xs T) T} と宣言していた。
scheme を具体化しても本体がそのまま返る。また \texttt{gen G T (Forall xs T)} に
側条件が無く、任意の変数を量化できた。両方が恒等写像なので、アルゴリズム的
関係と宣言的関係は構成子ごとに同一となり、健全性の定理は言い換えに過ぎなかった。

\paragraph{型安全性の対象言語に多相性が無かった。}
\texttt{HMTypeSafety.v} の型は \texttt{TInt} と \texttt{TArrow} だけで、型変数が
存在しなかった。ラムダは型注釈を持ち、\texttt{let} も単相であった。証明されて
いたのは、単純型付きラムダ計算に \texttt{let} と加算を足したものの型安全性である。

以下でどちらも修復する。

% =====================================================================
\section{型推論の健全性: 文法}
\label{sec:sound-syntax}

\subsection{文法}

\[
\begin{array}{lcll}
  e     & ::= & x \mid n \mid \lambda x.\,e \mid e_1\,e_2
                \mid \mathtt{let}\;x = e_1\;\mathtt{in}\;e_2
        & \text{式}\\[2pt]
  \tau  & ::= & \tint \mid \alpha \mid \tau_1 \to \tau_2
        & \text{型}\\[2pt]
  \sigma & ::= & \forall \bar\alpha.\,\tau
        & \text{型スキーム}\\[2pt]
  \Gamma & ::= & \cdot \mid \Gamma, x{:}\sigma
        & \text{環境}
\end{array}
\]

$x$ と $\alpha$ は自然数、$\bar\alpha$ はその列である。単相 scheme は
$\forall \varepsilon.\,\tau$ で、紛れがなければ $\tau$ と書く。Coq では次のとおり。

\begin{lstlisting}
Inductive ty := TInt | TVar : nat -> ty | TArrow : ty -> ty -> ty.
Inductive expr := EVar : nat -> expr | EInt : nat -> expr
                | ELam : nat -> expr -> expr
                | EApp : expr -> expr -> expr
                | ELet : nat -> expr -> expr -> expr.
Inductive scheme := Sch : list nat -> ty -> scheme.
Definition env := list (nat * scheme).
\end{lstlisting}

環境を関数ではなく連想リストにしたのは、下の $\ftv(\Gamma)$ を計算可能に
するためである。

\subsection{自由型変数}

\[
\begin{array}{lcl}
  \ftv(\tint) &=& \emptyset \\
  \ftv(\alpha) &=& \{\alpha\} \\
  \ftv(\tau_1 \to \tau_2) &=& \ftv(\tau_1) \cup \ftv(\tau_2) \\[2pt]
  \ftv(\forall\bar\alpha.\,\tau) &=& \ftv(\tau) \setminus \bar\alpha \\
  \ftv(\Gamma) &=& \textstyle\bigcup_{x:\sigma \,\in\, \Gamma} \ftv(\sigma)
\end{array}
\]

Coq では \texttt{ftv}、\texttt{ftv\_scheme}、\texttt{ftv\_env} である。差集合は
真偽値の所属判定 \texttt{memb} による \texttt{filter} で表し、両者を
\texttt{memb\_In} が結ぶ。

\subsection{具体化と一般化}

代入 $s$ は型変数を型へ写し、型に準同型に作用する。$\bar\alpha$ に属さない
すべての $\beta$ について $s(\beta) = \beta$ であるとき、$s$ は $\bar\alpha$ の
\emph{上でのみ作用する}という。

\begin{mathpar}
  \inferrule*[right=Inst]
    {s \text{ は } \bar\alpha \text{ の上でのみ作用する}}
    {\forall\bar\alpha.\,\tau \;\inst\; s(\tau)}
  \and
  \inferrule*[right=Gen]
    {\bar\alpha \cap \ftv(\Gamma) = \emptyset}
    {\gen(\Gamma, \tau) \ni \forall\bar\alpha.\,\tau}
\end{mathpar}

\textsc{Inst} が第一の欠陥の修復である。scheme の本体が実際に置換される。
ここから繰り返し使う帰結として、単相 scheme はただ一つの具体化しか持たない。
\[
  \forall\varepsilon.\,\tau \inst \tau'
  \quad\Longrightarrow\quad
  \tau' = \tau
\]
これが \texttt{inst\_mono\_inv} であり、証明は \texttt{only\_on\_nil}
（$\varepsilon$ の上でのみ作用する代入は恒等）と \texttt{appT\_id} から出る。

\textsc{Gen} が第二の欠陥の修復、すなわち Damas--Milner の側条件である。
空回りしていない。$\Gamma = \{x_0 : \alpha_7\}$ のとき
$\gen(\Gamma, \alpha_7) \ni \forall\alpha_7.\,\alpha_7$ は\emph{導出不能}であり、
これが \texttt{gen\_rejects\_env\_variable} である。以前の定義ではこれが導けた。

アルゴリズムは \textsc{Gen} が許す自由を持たず、最大一般化に確定する。
\[
  \genmax(\Gamma,\tau) \;=\; \forall\,(\ftv(\tau) \setminus \ftv(\Gamma)).\,\tau
\]

\section{型推論の健全性: 規則}

\subsection{宣言的体系}

\begin{mathpar}
  \inferrule*[right=T-Var]
    {\Gamma(x) = \sigma \\ \sigma \inst \tau}
    {\Gamma \vdash x : \tau}
  \and
  \inferrule*[right=T-Int]
    { }
    {\Gamma \vdash n : \tint}
  \and
  \inferrule*[right=T-Lam]
    {\Gamma, x{:}\tau_1 \vdash e : \tau_2}
    {\Gamma \vdash \lambda x.\,e : \tau_1 \to \tau_2}
  \and
  \inferrule*[right=T-App]
    {\Gamma \vdash e_1 : \tau_1 \to \tau_2 \\ \Gamma \vdash e_2 : \tau_1}
    {\Gamma \vdash e_1\,e_2 : \tau_2}
  \and
  \inferrule*[right=T-Let]
    {\Gamma \vdash e_1 : \tau_1 \\
     \gen(\Gamma,\tau_1) \ni \sigma \\
     \Gamma, x{:}\sigma \vdash e_2 : \tau_2}
    {\Gamma \vdash \mathtt{let}\;x = e_1\;\mathtt{in}\;e_2 : \tau_2}
\end{mathpar}

\subsection{アルゴリズム的体系}

\texttt{let} 以外は同じ規則である。\textsc{I-Let} には scheme を選ぶ自由が無い。

\begin{mathpar}
  \inferrule*[right=I-Let]
    {\Gamma \vdash_{\!I} e_1 : \tau_1 \\
     \Gamma, x{:}\genmax(\Gamma,\tau_1) \vdash_{\!I} e_2 : \tau_2}
    {\Gamma \vdash_{\!I} \mathtt{let}\;x = e_1\;\mathtt{in}\;e_2 : \tau_2}
\end{mathpar}

\textsc{T-Let} は \textsc{I-Let} が持たない前提を一つ余分に持つ。したがって
両者はもはや同じ規則ではなく、健全性には証明すべきことがある。

\section{型推論の健全性: 証明の進み方}
\label{sec:sound-proof}

\begin{theorem}[\texttt{hm\_infer\_sound}]
  $\Gamma \vdash_{\!I} e : \tau$ ならば $\Gamma \vdash e : \tau$ である。
\end{theorem}

\begin{proof}
\texttt{intros G e T H} のあと \texttt{induction H} で
$\Gamma \vdash_{\!I} e : \tau$ の導出に関する帰納法に入る。五つの場合を順に見る。

\medskip
\emph{第 1 場合 \textsc{I-Var}。} 文脈に $H_1 : \Gamma(x) = \sigma$ と
$H_2 : \sigma \inst \tau$ があり、目標は $\Gamma \vdash x : \tau$。
\textsc{T-Var} は同じ二つの前提を要求するので、\texttt{eapply TyVar; eauto} が
そのまま当てて閉じる。具体化の前提はアルゴリズム側と宣言側で共通なので、
書き換えは要らない。

\medskip
\emph{第 2 場合 \textsc{I-Int}。} 目標 $\Gamma \vdash n : \tint$ は前提を持たない
\textsc{T-Int} の結論そのもので、\texttt{constructor} で閉じる。

\medskip
\emph{第 3 場合 \textsc{I-Lam}。} 帰納法の仮定は
$\mathit{IH} : \Gamma, x{:}\forall\varepsilon.\tau_1 \vdash e : \tau_2$。
\textsc{T-Lam} を \texttt{apply} すると目標がちょうど $\mathit{IH}$ になる。

\medskip
\emph{第 4 場合 \textsc{I-App}。} 帰納法の仮定二つを \textsc{T-App} の二つの前提に
当てる。中間の型は \texttt{eapply} がメタ変数として置き、第一の仮定を当てた時点で
確定する。

\medskip
\emph{第 5 場合 \textsc{I-Let}。ここが本題である。} 手元にあるのは
\[
  \mathit{IH}_1 : \Gamma \vdash e_1 : \tau_1,
  \qquad
  \mathit{IH}_2 : \Gamma, x{:}\genmax(\Gamma,\tau_1) \vdash e_2 : \tau_2
\]
の二つだけで、$\gen$ についての前提は\emph{無い}。アルゴリズム側の \textsc{I-Let} が
それを要求しないからである。一方 \textsc{T-Let} は三つの前提を持ち、その中央が
\[
  \gen(\Gamma,\tau_1) \ni \sigma
\]
である。\texttt{eapply TyLet} すると $?\sigma$ はメタ変数として置かれるが、
第三の前提に $\mathit{IH}_2$ を当てた時点で
$?\sigma := \genmax(\Gamma,\tau_1)$ と確定する。したがって残る目標は
\[
  \gen(\Gamma,\tau_1) \ni \genmax(\Gamma,\tau_1)
\]
すなわち\emph{アルゴリズムが計算した scheme が、宣言的体系の側条件を満たしているか}
である。これは帰納法の仮定からは出ない。補題 \texttt{gen\_max\_ok} が要る。

\medskip
証明木で書くと、この場合が何をしているかが一目で分かる。左が手元にある
アルゴリズム側の導出、右が組み立てる宣言側の導出である。二重線は帰納法の仮定で
移した部分、太字が新たに証明する部分を表す。

\begin{mathpar}
  \inferrule*[right=I-Let]
    {\Gamma \vdash_{\!I} e_1 : \tau_1 \\
     \Gamma, x{:}\genmax(\Gamma,\tau_1) \vdash_{\!I} e_2 : \tau_2}
    {\Gamma \vdash_{\!I} \mathtt{let}\;x = e_1\;\mathtt{in}\;e_2 : \tau_2}
  \and
  {}\rightsquigarrow{}
  \and
  \inferrule*[right=T-Let]
    {\inferrule*[Right={\scriptsize $\mathit{IH}_1$}]{ }{\Gamma \vdash e_1 : \tau_1} \\
     \inferrule*[Right={\scriptsize \textbf{gen\_max\_ok}}]{ }
       {\gen(\Gamma,\tau_1) \ni \genmax(\Gamma,\tau_1)} \\
     \inferrule*[Right={\scriptsize $\mathit{IH}_2$}]{ }
       {\Gamma, x{:}\genmax(\Gamma,\tau_1) \vdash e_2 : \tau_2}}
    {\Gamma \vdash \mathtt{let}\;x = e_1\;\mathtt{in}\;e_2 : \tau_2}
\end{mathpar}

左の導出には前提が二つしかないのに、右は三つ要る。その差が
\texttt{gen\_max\_ok} である。

\medskip
以前の版では \textsc{I-Let} も \textsc{T-Let} と同じ $\gen$ の前提を持っており、
両者の木は形まで同一であった。だから右の木は左の木のラベルを付け替えるだけで作れ、
証明に義務が生じなかった。アルゴリズム側から選択の自由を取り上げたことで、初めて
ここに埋めるべき穴が現れる。
\end{proof}

\begin{lemma}[\texttt{gen\_max\_ok}]
  すべての $\Gamma, \tau$ について $\gen(\Gamma,\tau) \ni \genmax(\Gamma,\tau)$。
\end{lemma}

\begin{proof}
\texttt{unfold gen\_max} で量化される列が
$\mathtt{filter}\,(\lambda a.\,\neg\,\mathtt{memb}\;a\;\ftv(\Gamma))\;\ftv(\tau)$
であることを出し、\textsc{Gen} を \texttt{apply} する。残る目標は、その列の任意の
要素 $a$ について $a \notin \ftv(\Gamma)$ を示すことである。

$a$ が列に属するという仮定に \texttt{filter\_In} を適用すると、二つの情報に分かれる。
$a \in \ftv(\tau)$ であること（使わない）と、filter の述語が $a$ について
\texttt{true} を返すこと、すなわち
\[
  \neg\,(\mathtt{memb}\;a\;\ftv(\Gamma)) = \mathtt{true}
\]
である。これに \texttt{negb\_true\_iff} を適用して
$\mathtt{memb}\;a\;\ftv(\Gamma) = \mathtt{false}$ を得る。最後に
\texttt{memb\_false\_notin}（真偽値の所属判定が \texttt{false} なら
命題としての所属も成り立たない、という補題で、\texttt{memb\_In} から出る）を
適用すれば $a \notin \ftv(\Gamma)$ が得られる。これがちょうど \textsc{Gen} の
側条件である。

真偽値としての判定と命題としての所属を行き来する二段（\texttt{negb\_true\_iff} と
\texttt{memb\_false\_notin}）が必要になるのは、$\ftv$ を計算可能にするために
環境を連想リストにし、所属を真偽値で書いたことの代償である。
\end{proof}

\subsection{実例: 多相性が実際に使われる}

scheme $\forall\alpha.\,\alpha \to \alpha$ は、$\alpha_0$ の上でのみ作用する
代入 $s_1 = [\alpha_0 := \tint]$ と $s_2 = [\alpha_0 := \tint \to \tint]$ から、
二つの具体化を得る。
\[
  \forall\alpha_0.\,\alpha_0 \to \alpha_0 \;\inst\; \tint \to \tint,
  \qquad
  \forall\alpha_0.\,\alpha_0 \to \alpha_0 \;\inst\;
  (\tint \to \tint) \to (\tint \to \tint)
\]
これを用いて、プログラム
\[
  \mathtt{let}\;\mathit{id} = \lambda x.\,x\;\mathtt{in}\;
  (\mathit{id}\;\mathit{id})\;42
\]
をアルゴリズム的体系で型 $\tint$ に導き、定理で宣言的体系へ移す。左の
$\mathit{id}$ は $(\tint \to \tint) \to (\tint \to \tint)$ で、右の
$\mathit{id}$ は $\tint \to \tint$ で使われる。補題
\texttt{selfapp\_needs\_polymorphism} は、単相束縛ではこれが不可能であることを
示す。\texttt{inst\_mono\_inv} により二つの出現は同じ型を受け取らねばならず、
その二つの型は異なるからである。

% =====================================================================
\section{型安全性: 文法}
\label{sec:safe-syntax}

\subsection{文法}

\[
\begin{array}{lcll}
  t & ::= & x \mid n \mid t_1 + t_2 \mid \lambda x.\,t \mid t_1\,t_2
            \mid \mathtt{let}\;x = t_1\;\mathtt{in}\;t_2
    & \text{項}\\[2pt]
  v & ::= & n \mid \lambda x.\,t
    & \text{値}\\[2pt]
  \tau & ::= & \tint \mid \alpha \mid \#i \mid \tau_1 \to \tau_2
    & \text{型}\\[2pt]
  \sigma & ::= & \forall^{n}.\,\tau
    & \text{スキーム}
\end{array}
\]

ラムダは型注釈を持たない。言語は型無しで、関係がそれに型を付ける。実際の
HM 体系と同じ形である。

型変数は\emph{二種類}ある。$\alpha$ は自由型変数（\texttt{TFree}）、$\#i$ は
囲む scheme が束縛する $i$ 番目の変数（\texttt{TBound}）、すなわち de Bruijn
指標である。scheme $\forall^{n}.\,\tau$ は $\tau$ の中の $\#0,\dots,\#(n-1)$ を
束縛する。

この二分割が、開発を短く保っている仕掛けである。具体化は $\#i$ しか置き換えず、
置き換えて入る型には $\alpha$ しか現れないので、自由変数が量化子に捕獲される
ことが原理的に起こらない。新鮮性の側条件も型代入補題も、どこにも必要ない。

\subsection{具体化}

$\bar\tau = \tau_0 \dots \tau_{n-1}$ に対して次を定める。
\[
\begin{array}{lcl}
  \mathtt{open}\,\bar\tau\,\tint &=& \tint \\
  \mathtt{open}\,\bar\tau\,\alpha &=& \alpha \\
  \mathtt{open}\,\bar\tau\,\#i &=& \tau_i \quad (i \geq n \text{ なら } \#i) \\
  \mathtt{open}\,\bar\tau\,(\tau_1 \to \tau_2) &=&
    \mathtt{open}\,\bar\tau\,\tau_1 \to \mathtt{open}\,\bar\tau\,\tau_2
\end{array}
\]

\begin{mathpar}
  \inferrule*[right=Inst]
    {|\bar\tau| = n}
    {\forall^{n}.\,\tau \;\inst\; \mathtt{open}\,\bar\tau\,\tau}
\end{mathpar}

$\mathtt{open}\,\varepsilon\,\tau = \tau$ が無条件に成り立つ
（\texttt{open\_nil}）ので、単相 scheme $\forall^{0}.\,\tau$ はやはり $\tau$
ただ一つの具体化を持つ（\texttt{inst\_mono}、\texttt{inst\_mono\_inv}）。

\section{型安全性: 規則}

\subsection{型付け}

\begin{mathpar}
  \inferrule*[right=S-Var]
    {\Gamma(x) = \sigma \\ \sigma \inst \tau}
    {\Gamma \vdash x : \tau}
  \and
  \inferrule*[right=S-Int]
    { }
    {\Gamma \vdash n : \tint}
  \and
  \inferrule*[right=S-Add]
    {\Gamma \vdash t_1 : \tint \\ \Gamma \vdash t_2 : \tint}
    {\Gamma \vdash t_1 + t_2 : \tint}
  \and
  \inferrule*[right=S-Lam]
    {\Gamma, x{:}\forall^{0}.\,\tau_1 \vdash t : \tau_2}
    {\Gamma \vdash \lambda x.\,t : \tau_1 \to \tau_2}
  \and
  \inferrule*[right=S-App]
    {\Gamma \vdash t_1 : \tau_1 \to \tau_2 \\ \Gamma \vdash t_2 : \tau_1}
    {\Gamma \vdash t_1\,t_2 : \tau_2}
  \and
  \inferrule*[right=S-Let]
    {\forall \bar\tau.\; |\bar\tau| = n \Rightarrow
       \Gamma \vdash t_1 : \mathtt{open}\,\bar\tau\,\tau_0 \\\\
     \Gamma, x{:}\forall^{n}.\,\tau_0 \vdash t_2 : \tau}
    {\Gamma \vdash \mathtt{let}\;x = t_1\;\mathtt{in}\;t_2 : \tau}
\end{mathpar}

\textsc{S-Let} は一般化を意味論的に読んだものである。$\ftv(\Gamma)$ から
scheme を計算するのではなく、$t_1$ が scheme の\emph{全ての}具体化で型付く
ことを要求し、$t_2$ の中の $x$ の各出現が必要な具体化を選べるようにする。
この前提は全称量化された仮定だが、Coq はこれを強正値の出現として受理し、
期待どおりの帰納法の仮定を生成する。

\subsection{操作意味論}

値呼び、左から右。

\begin{mathpar}
  \inferrule*[right=E-Add]
    { }
    {n + m \step n{+}m}
  \and
  \inferrule*[right=E-Add1]
    {t_1 \step t_1'}
    {t_1 + t_2 \step t_1' + t_2}
  \and
  \inferrule*[right=E-Add2]
    {t \step t'}
    {v + t \step v + t'}
  \and
  \inferrule*[right=E-Beta]
    { }
    {(\lambda x.\,t)\,v \step \subst{t}{x}{v}}
  \and
  \inferrule*[right=E-App1]
    {t_1 \step t_1'}
    {t_1\,t_2 \step t_1'\,t_2}
  \and
  \inferrule*[right=E-App2]
    {t \step t'}
    {v\,t \step v\,t'}
  \and
  \inferrule*[right=E-Let]
    { }
    {\mathtt{let}\;x = v\;\mathtt{in}\;t \step \subst{t}{x}{v}}
  \and
  \inferrule*[right=E-Let1]
    {t_1 \step t_1'}
    {\mathtt{let}\;x = t_1\;\mathtt{in}\;t_2 \step
     \mathtt{let}\;x = t_1'\;\mathtt{in}\;t_2}
\end{mathpar}

代入 $\subst{t}{x}{v}$ は、$\lambda x$ と $\mathtt{let}\;x = \cdot$ の本体で
止まる、隠蔽による捕獲回避の通常の定義である。

\section{型安全性: 証明の進み方}
\label{sec:safe-proof}

\subsection{環境についての補題}

環境は $\mathbb{N} \to \sigma_\bot$ の関数なので、必要な四つの事実は一行ずつで
済む。\texttt{extend\_eq}、\texttt{extend\_neq}、\texttt{extend\_shadow}
$(\Gamma, x{:}\sigma_1, x{:}\sigma_2 = \Gamma, x{:}\sigma_2)$、
\texttt{extend\_permute} $(x \neq y \Rightarrow
\Gamma, x{:}\sigma_x, y{:}\sigma_y = \Gamma, y{:}\sigma_y, x{:}\sigma_x)$ で、
いずれも $\mathtt{Nat.eqb}$ の場合分けで終わる。その上に
\texttt{context\_invariance}（各点で等しい環境は同じ項に同じ型を与える）と
\texttt{weakening} が乗り、どちらも型付け導出に関する帰納法である。

\subsection{代入補題}

\begin{lemma}[\texttt{substitution\_preserves\_typing}]
  $\Gamma, x{:}\sigma \vdash t : \tau$ であり、かつ $\sigma \inst \tau'$ なる
  \emph{すべての} $\tau'$ について $\vdash v : \tau'$ であるなら、
  $\Gamma \vdash \subst{t}{x}{v} : \tau$ である。
\end{lemma}

$v$ への仮定が、多相束縛でこの補題を働かせている。$x$ が使われうる各具体化で
$v$ が型付いていなければならない。

\begin{proof}
$\tau$ と $\Gamma$ を \texttt{generalize dependent} してから、導出ではなく\emph{項}
$t$ に関する帰納法に入る。項に関する帰納法を取るのは、各場合が $v$ への仮定を
\emph{別の型で}使い直せるようにするためである。導出に関する帰納法だと、$v$ の型が
一つに固定されてしまい多相の場合が通らない。各場合ではまず \texttt{simpl} で
\texttt{subst} を一段展開し、\texttt{inversion Ht} で型付け導出を分解する。

\medskip
\emph{第 1 場合 $t = y$（変数）。} \texttt{inversion} から
$H_l : (\Gamma, x{:}\sigma)(y) = \mathtt{Some}\;\sigma_0$ と
$H_i : \sigma_0 \inst \tau$ が出る。$\mathtt{Nat.eqb}\;x\;y$ で場合分けする。

\texttt{true} の枝では $x = y$ なので、$H_l$ に \texttt{extend\_eq} を書き換えると
$\mathtt{Some}\;\sigma = \mathtt{Some}\;\sigma_0$ となり、\texttt{inversion} で
$\sigma_0 = \sigma$。すると $H_i$ は $\sigma \inst \tau$ になる。$v$ への仮定は
「$\sigma$ の\emph{全ての}具体化 $\tau'$ について $\vdash v : \tau'$」であったから、
$H_i$ を渡して $\vdash v : \tau$ を得る。あとは \texttt{weaken\_empty} で $\Gamma$ へ
持ち上げる。仮定が「全ての具体化について」の形であることが、ここで効いている。

\texttt{false} の枝では $x \neq y$。\texttt{extend\_neq} で束縛を捨て、
\textsc{S-Var} を $H_l$ と $H_i$ で組み直す。

\medskip
\emph{第 2 場合 $t = n$。} \textsc{S-Int} が前提なしに与える。

\medskip
\emph{第 3 場合 $t = t_1 + t_2$。} \textsc{S-Add} を組み直し、両辺に帰納法の仮定を
$v$ への仮定とともに当てる。

\medskip
\emph{第 4 場合 $t = \lambda y.\,t'$。} \texttt{inversion} から
$H_b : (\Gamma, x{:}\sigma), y{:}\forall^{0}.\tau_1 \vdash t' : \tau_2$。
$\mathtt{Nat.eqb}\;x\;y$ で場合分けする。

$y = x$ のとき、\texttt{subst} の定義により代入はラムダの内側へ入らない。目標は
$\Gamma \vdash \lambda x.\,t' : \tau_1 \to \tau_2$ で、\textsc{S-Lam} の後に残るのは
$\Gamma, x{:}\forall^{0}.\tau_1 \vdash t' : \tau_2$。手元の $H_b$ は
$(\Gamma, x{:}\sigma), x{:}\forall^{0}.\tau_1$ の下だが、\texttt{extend\_shadow} が
\[
  (\Gamma, x{:}\sigma), x{:}\forall^{0}.\tau_1 \;=\; \Gamma, x{:}\forall^{0}.\tau_1
\]
を各点で与えるので、\texttt{context\_invariance} で環境を移せる。

$y \neq x$ のとき、代入は内側へ入る。帰納法の仮定は環境が
$(\Gamma, y{:}\forall^{0}.\tau_1), x{:}\sigma$ の形であることを求めるが、$H_b$ は
順序が逆である。\texttt{extend\_permute}（副条件 $x \neq y$）で入れ替えてから
仮定を当てる。二つの枝を証明木で並べると次のようになる。

\begin{mathpar}
  \inferrule*[right=S-Lam]
    {\inferrule*[Right={\scriptsize context\_invariance $+$ extend\_shadow}]
       {(\Gamma, x{:}\sigma), x{:}\forall^{0}.\tau_1 \vdash t' : \tau_2}
       {\Gamma, x{:}\forall^{0}.\tau_1 \vdash t' : \tau_2}}
    {\Gamma \vdash \lambda x.\,t' : \tau_1 \to \tau_2}
  \and
  \inferrule*[right=S-Lam]
    {\inferrule*[Right={\scriptsize $\mathit{IH}$}]
       {\inferrule*[Right={\scriptsize context\_invariance $+$ extend\_permute}]
          {(\Gamma, x{:}\sigma), y{:}\forall^{0}.\tau_1 \vdash t' : \tau_2}
          {(\Gamma, y{:}\forall^{0}.\tau_1), x{:}\sigma \vdash t' : \tau_2}}
       {\Gamma, y{:}\forall^{0}.\tau_1 \vdash \subst{t'}{x}{v} : \tau_2}}
    {\Gamma \vdash \lambda y.\,\subst{t'}{x}{v} : \tau_1 \to \tau_2}
\end{mathpar}

左が $y = x$（代入が止まる）、右が $y \neq x$（代入が入る）である。どちらも
\textsc{S-Lam} の下で環境を一段付け替えており、違いは付け替えに使う補題が
\texttt{extend\_shadow} か \texttt{extend\_permute} か、そして帰納法の仮定を
挟むかどうかだけである。

\medskip
\emph{第 5 場合 $t = t_1\,t_2$。} \textsc{S-App} を \texttt{eapply} して両辺に
帰納法の仮定を当てる。

\medskip
\emph{第 6 場合 $t = \mathtt{let}\;y = t_1\;\mathtt{in}\;t_2$。} \textsc{S-Let} を
\texttt{eapply} すると二つの目標に分かれる。

第一の目標は「代入後の $t_1$ が scheme の\emph{全ての}具体化で型付く」ことである。
$\bar\tau$ を導入し、\textsc{S-Let} の全称量化された前提をその $\bar\tau$ で具体化して
$(\Gamma, x{:}\sigma) \vdash t_1 : \mathtt{open}\,\bar\tau\,\tau_0$ を取り出し、
$t_1$ の帰納法の仮定に渡す。前提が全称量化されているので、$\bar\tau$ ごとに一度ずつ
帰納法の仮定を使い直すことになる。

第二の目標は本体で、第 4 場合と同じ二分岐である。$y = x$ なら
\texttt{extend\_shadow}、$y \neq x$ なら \texttt{extend\_permute} を
\texttt{context\_invariance} と組み合わせる。

\medskip
全体を通して働いているのは、\emph{束縛が名前を隠すか順序を入れ替えられるか}という
二つの事実と、\emph{$v$ への仮定が全ての具体化について述べられている}ことである。
後者が多相 \texttt{let} を通している。
\end{proof}

\subsection{標準形の補題}

$\mathtt{value}\,v$ と型付け導出の両方を inversion して得る。$\tint$ 型の値は
整数リテラル、$\tau_1 \to \tau_2$ 型の値はラムダである。他の組み合わせは、
\textsc{S-Int} と \textsc{S-Lam} が異なる型構成子を作ることから反駁される。

\subsection{Progress}

\begin{theorem}[\texttt{progress}]
  $\vdash t : \tau$ ならば、$t$ は値であるか、ある $t'$ へ $t \step t'$ と
  進む。
\end{theorem}

\begin{proof}
\texttt{remember empty as G} で環境を $\emptyset$ に固定してから型付け導出に関する
帰納法に入る。この \texttt{remember} が要るのは、\textsc{S-Lam} と \textsc{S-Let} の
部分導出が拡張された環境の下にあり、そのままでは帰納法の仮定を空環境について
使えないからである。六つの場合を順に見る。

\medskip
\emph{\textsc{S-Var}。} 前提は $\emptyset(x) = \mathtt{Some}\;\sigma$ だが、
$\emptyset$ は定義上すべての $x$ に $\mathtt{None}$ を返す。構成子が違うので
\texttt{discriminate} でこの場合ごと消える。閉じた項に自由変数は無い、という事実が
ここに現れる。

\medskip
\emph{\textsc{S-Int}、\textsc{S-Lam}。} それぞれ \textsc{V-Int}、\textsc{V-Lam} に
より値。左の選択肢を取る。

\medskip
\emph{\textsc{S-Add}。} 右の選択肢を主張し、$t_1$ の帰納法の仮定で場合分けする。
$t_1$ が進むなら \textsc{E-Add1}。$t_1$ が値なら $t_2$ の仮定で分け、$t_2$ が進むなら
\textsc{E-Add2}（$t_1$ が値であることがその前提として要る）。両方値なら
\texttt{canonical\_int} を二度使って $t_1 = n$、$t_2 = m$ を取り出し、
\textsc{E-Add} が発火する。証人は $\mathtt{TmInt}\,(n+m)$。

\medskip
\emph{\textsc{S-App}。} 同じ運びで、両方値のときに \texttt{canonical\_arrow} を
$t_1$ に適用して $t_1 = \lambda x.\,t$ を取り出す。これで \textsc{E-Beta} が使え、
証人は $\subst{t}{x}{t_2}$。ここが標準形の補題を要する箇所で、「型が付いているのに
関数の位置に関数でないものが来る」という詰まり方が起こらないことを保証している。

\medskip
\emph{\textsc{S-Let}。ここだけ手当てが要る。} $t_1$ についての帰納法の仮定は、
\textsc{S-Let} の前提が全称量化されているために、それ自身
\[
  \forall \bar\tau.\; |\bar\tau| = n \Rightarrow
  \bigl(\text{$t_1$ は値であるか一歩進む}\bigr)
\]
という形をしている。使う前に $\bar\tau$ を一つ与えねばならない。長さが $n$ でさえ
あれば結論は $\bar\tau$ に依らないので、どれでもよい。証明は
$\mathtt{repeat}\;\tint\;n$ を渡し、その長さが $n$ であることを
\texttt{repeat\_length} で示す。これで仮定が使える形になり、$t_1$ が値なら
\textsc{E-Let}、進むなら \textsc{E-Let1} が働く。

\medskip
どの場合でも「値である」か「一歩進む証人を作る」かのいずれかが得られた。
\end{proof}

\subsection{Preservation}

\begin{theorem}[\texttt{preservation}]
  $\vdash t : \tau$ かつ $t \step t'$ ならば $\vdash t' : \tau$ である。
\end{theorem}

\begin{proof}
$\tau$ を \texttt{generalize dependent} してから step の導出に関する帰納法に入り、
各場合でまず型付け導出を \texttt{inversion} する。step の構成子は八つ。

\medskip
\emph{\textsc{E-Add}。} \texttt{inversion} から $\tau = \tint$。目標
$\vdash n{+}m : \tint$ は \textsc{S-Int} が与える。

\medskip
\emph{\textsc{E-Add1}、\textsc{E-Add2}。} \textsc{S-Add} を組み直し、進む側に
帰納法の仮定を当て、他方はそのまま持ち越す。

\medskip
\emph{\textsc{E-Beta}。} \textsc{S-App} の \texttt{inversion} で
$\vdash \lambda x.\,t : \tau_1 \to \tau$ と $\vdash v : \tau_1$ を得、続いて前者を
\textsc{S-Lam} で \texttt{inversion} して
$\emptyset, x{:}\forall^{0}.\tau_1 \vdash t : \tau$ を取り出す。代入補題は
「$v$ が $\forall^{0}.\tau_1$ の\emph{全ての}具体化で型付く」ことを求めるが、
\texttt{inst\_mono\_inv} により具体化は $\tau_1$ ただ一つなので、手元の
$\vdash v : \tau_1$ で足りる。ラムダ束縛が単相であることが、ここで全称量化された
仮定を一つの型へ潰している。

\begin{mathpar}
  \inferrule*[right={\scriptsize 代入補題}]
    {\inferrule*[Right={\scriptsize S-Lam の inversion}]
       {\inferrule*[Right={\scriptsize S-App の inversion}]
          {\vdash (\lambda x.\,t)\,v : \tau}
          {\vdash \lambda x.\,t : \tau_1 \to \tau}}
       {\emptyset, x{:}\forall^{0}.\tau_1 \vdash t : \tau}
     \\
     \inferrule*[Right={\scriptsize inst\_mono\_inv}]
       {\vdash v : \tau_1}
       {\forall \tau'.\; \forall^{0}.\tau_1 \inst \tau' \Rightarrow \vdash v : \tau'}}
    {\vdash \subst{t}{x}{v} : \tau}
\end{mathpar}

\medskip
\emph{\textsc{E-App1}、\textsc{E-App2}。} \textsc{S-App} を組み直す。中間の型は
\texttt{eapply} がメタ変数として扱い、他方の前提を当てた時点で確定する。

\medskip
\emph{\textsc{E-Let}。ここが多相の場合である。} \textsc{S-Let} を
\texttt{inversion} すると、
\[
  \forall \bar\tau.\; |\bar\tau| = n \Rightarrow
    \vdash v : \mathtt{open}\,\bar\tau\,\tau_0
  \qquad\text{と}\qquad
  \emptyset, x{:}\forall^{n}.\tau_0 \vdash t_2 : \tau
\]
が同時に得られる。代入補題が求めるのは「$v$ が $\forall^{n}.\tau_0$ の全ての具体化で
型付く」ことであり、$\forall^{n}.\tau_0 \inst \tau'$ を \texttt{inversion} すると
$\tau' = \mathtt{open}\,\bar\tau\,\tau_0$ で $|\bar\tau| = n$ という形が出るので、
第一の前提をその $\bar\tau$ で具体化すればそのまま渡せる。\textbf{追加の作業が要らない。}

\begin{mathpar}
  \inferrule*[right={\scriptsize 代入補題}]
    {\inferrule*[Right={\scriptsize S-Let の inversion（第二前提）}]
       {\vdash \mathtt{let}\;x = v\;\mathtt{in}\;t_2 : \tau}
       {\emptyset, x{:}\forall^{n}.\tau_0 \vdash t_2 : \tau}
     \\
     \inferrule*[Right={\scriptsize S-Let の inversion（第一前提）}]
       {\vdash \mathtt{let}\;x = v\;\mathtt{in}\;t_2 : \tau}
       {\forall \bar\tau.\;|\bar\tau| = n \Rightarrow
          \vdash v : \mathtt{open}\,\bar\tau\,\tau_0}}
    {\vdash \subst{t_2}{x}{v} : \tau}
\end{mathpar}

\textsc{E-Beta} の木と見比べると違いがはっきりする。あちらは \texttt{inversion} を
二段重ね、さらに \texttt{inst\_mono\_inv} で全称量化を一つの型へ潰す必要があった。
こちらは一段の \texttt{inversion} で二つの前提が出そろい、第一前提が既に
「全ての具体化について」の形をしている。
\textsc{S-Let} を「全ての具体化で型付く」という意味論的な形で述べたことの見返りが
ここにある。構文的な \textsc{Gen} 規則を採ると、この箇所で型代入補題と新鮮性の議論が
必要になる。

\medskip
\emph{\textsc{E-Let1}。} \textsc{S-Let} を組み直す。第一の前提は、古い前提を各
$\bar\tau$ で具体化し帰納法の仮定を通して作る。本体の前提は $\tau_0$ も $n$ も
変わらないのでそのまま持ち越す。

\medskip
八つの場合が閉じる。合同規則の六つは規則を組み直すだけで、実質的な内容は
\textsc{E-Beta} と \textsc{E-Let} の二つ、いずれも代入補題への帰着である。
\end{proof}

\subsection{実例: 多相性が実際に使われる}

$\mathit{id} = \lambda x.\,x$ が $\forall^{1}.\,\#0 \to \#0$ の全ての具体化で
型付くことを示す。任意の $\bar\tau = \tau$ について
$\mathtt{open}\,[\tau]\,(\#0 \to \#0) = \tau \to \tau$ であり、\textsc{S-Lam} に
続いて \texttt{inst\_mono} つきの \textsc{S-Var} を使えば
$\vdash \lambda x.\,x : \tau \to \tau$ が一般に導ける。この前提が
\textsc{S-Let} を養い、
\[
  \vdash \mathtt{let}\;\mathit{id} = \lambda x.\,x\;\mathtt{in}\;
  (\mathit{id}\;\mathit{id})\;42 \;:\; \tint
\]
を与える。二つの出現はそれぞれ $[\tint \to \tint]$ と $[\tint]$ で具体化される。
前の版の単相体系ではこれは型付かない。最後に \textsc{E-Let} で一歩進め、
\texttt{preservation} で型を付け直している。

% =====================================================================
\section{一般化の二つの見方}

二つのファイルは一般化を意図的に別の角度から扱っている。

\begin{center}
\begin{tabular}{lll}
  & \texttt{HMSoundness.v} & \texttt{HMTypeSafety.v} \\
  \hline
  形 & 構文的 & 意味論的 \\
  条件 & $\bar\alpha \cap \ftv(\Gamma) = \emptyset$
       & 全ての具体化で型付く \\
  量化変数の束縛 & 名前付き & de Bruijn 指標 \\
  環境 & 連想リスト & 関数 \\
  $\ftv(\Gamma)$ の要否 & 必要 & 不要 \\
\end{tabular}
\end{center}

意味論的な形は、型代入補題と新鮮性の機構なしに安全性の証明を閉じるための
選択である。標準的な定式化ではあるが、構文的な \textsc{Gen} 規則そのものでは
ない。その規則を求める読者は \texttt{HMSoundness.v} を見られたい。両者を
繋ぐこと、すなわち構文的な一般化が意味論的な条件を満たすことの証明は、自然な
次の一歩であり、ここでは行っていない。

\section{限界の突破: 再帰、多段の安全性、occurs check}
\label{sec:beyond}

前節までの体系には、繰り返し「対象外」と記してきた欠落があった。本節はそのうち
三つを埋める。いずれも Coq で証明済みで、追加公理に依存しない。

\subsection{再帰を入れる}

MiniML の中核は \texttt{let rec} であるのに、型安全性の証明に再帰が無かった。
再帰関数を項に加える。
\[
  t \;::=\; \dots \mid \mathtt{fix}\;f\;x.\,t
\]
$\mathtt{fix}\;f\;x.\,t$ は値であり、自分自身の名前 $f$ が本体で使える。

\begin{mathpar}
  \inferrule*[right=S-Fix]
    {\Gamma, f{:}\forall^{0}.(\tau_1 \to \tau_2), x{:}\forall^{0}.\tau_1
       \vdash t : \tau_2}
    {\Gamma \vdash \mathtt{fix}\;f\;x.\,t : \tau_1 \to \tau_2}
  \and
  \inferrule*[right=E-FixApp]
    {\mathrm{value}(v)}
    {(\mathtt{fix}\;f\;x.\,t)\,v \step
       \subst{\bigl(\subst{t}{x}{v}\bigr)}{f}{\mathtt{fix}\;f\;x.\,t}}
\end{mathpar}

\textsc{S-Fix} は、これから作る矢印型で $f$ を単相に束縛したうえで本体を型付ける。
\textsc{E-FixApp} は引数を代入し、続いて再帰を一段ほどく。

\paragraph{何が変わったか。}
\textsc{S-Fix} が増えたことで、既存の証明すべてに新しい場合が生じた。

\emph{標準形の補題。} 関数型の値はもはやラムダとは限らない。結論を選言に変える。
\[
  \mathrm{value}(v),\; \vdash v : \tau_1 \to \tau_2
  \;\Longrightarrow\;
  (\exists x, t.\; v = \lambda x.\,t) \;\lor\;
  (\exists f, x, t.\; v = \mathtt{fix}\;f\;x.\,t)
\]

\emph{progress。} 適用の場合がこの選言で分岐する。関数がラムダなら
\textsc{E-Beta}、再帰関数なら \textsc{E-FixApp} が発火する。証人が二通りになった
だけで、詰まらないことは変わらない。

\emph{代入補題。} $\mathtt{fix}\;f\;x.\,t$ は束縛を\emph{二つ}持つので、代入する
変数がどちらに一致するかで四つの場合に分かれる。$f$ と一致、$x$ と一致、両方と
一致、どちらとも異なる、である。前三者では代入が止まり、環境の付け替えを
\texttt{context\_invariance} で行う。最後の場合だけ代入が本体に入り、帰納法の
仮定を使う。

\emph{preservation。} \textsc{E-FixApp} の場合が新しい。代入補題を\emph{二度}
使う。まず外側の束縛 $x$ を剥がして引数 $v$ を代入し、次に残った束縛 $f$ を
剥がして再帰関数自身を代入する。環境
$(\Gamma, f{:}\tau_1 \to \tau_2), x{:}\tau_1$ の $x$ が外側にあるので、この順序
なら束縛を順に剥がすだけで済み、$f$ と $x$ の相異を仮定する必要がない。

\begin{mathpar}
  \inferrule*[right={\scriptsize 代入補題（$f$ を剥がす）}]
    {\inferrule*[Right={\scriptsize 代入補題（$x$ を剥がす）}]
       {(\emptyset, f{:}\tau_1 \to \tau_2), x{:}\tau_1 \vdash t : \tau_2
        \\ \vdash v : \tau_1}
       {\emptyset, f{:}\tau_1 \to \tau_2 \vdash \subst{t}{x}{v} : \tau_2}
     \\
     \vdash \mathtt{fix}\;f\;x.\,t : \tau_1 \to \tau_2}
    {\vdash \subst{\bigl(\subst{t}{x}{v}\bigr)}{f}{\mathtt{fix}\;f\;x.\,t}
       : \tau_2}
\end{mathpar}

\subsection{多段の安全性}

単段の preservation は「一歩進んでも型が保たれる」としか言わない。本来主張したい
のは「型の付いた閉じた項は、何歩進んでも\emph{詰まった状態}に到達しない」である。
\[
  \mathrm{stuck}(t) \;\equiv\;
    \neg\,\mathrm{value}(t) \;\land\; \neg\,\exists t'.\; t \step t'
\]

\begin{theorem}[\texttt{soundness}]
  $\vdash t : \tau$ かつ $t \step^{*} t'$ ならば $\neg\,\mathrm{stuck}(t')$。
\end{theorem}

\begin{proof}
まず多段版の preservation を、$t \step^{*} t'$ の導出に関する帰納法で示す。
反射の場合は仮定そのもの、推移の場合は単段の preservation を一度使ってから
帰納法の仮定に渡す。これで $\vdash t' : \tau$ が得られる。あとは $t'$ に
progress を適用すれば、$t'$ は値であるか一歩進むかのいずれかであり、どちらも
$\mathrm{stuck}(t')$ の二つの連言肢と矛盾する。
\end{proof}

これが「型の付いたプログラムは誤らない」の完全な形である。

\paragraph{停止性は主張していない。}
再帰が入った以上、型が付いても止まらない項が書ける。実際
$\mathtt{fix}\;f\;x.\,f\,x$ は任意の $\tau_1 \to \tau_2$ に型が付き、
$(\mathtt{fix}\;f\;x.\,f\,x)\,0$ は自分自身へ簡約される。この項について
\texttt{loop\_never\_stuck} を証明してある。永遠に走り続けることと詰まることは
別であり、型安全性が保証するのは後者が起きないことだけである。この区別が
はっきり見えるようになったのは、再帰を入れたからである。

\subsection{occurs check}

単一化は、変数 $\alpha$ がその型の内部に真に現れるとき、方程式
$\alpha = \tau$ を拒否しなければならない。これは約束事ではなく定理である。

\begin{theorem}[\texttt{occurs\_check\_sound}]
  $\alpha$ が $\tau_1 \to \tau_2$ に出現するならば、
  $s(\alpha) = s(\tau_1 \to \tau_2)$ を満たす代入 $s$ は存在しない。
\end{theorem}

\begin{proof}
型の大きさ $\mathrm{size}$ を、$\tint$ と変数を $1$、
$\mathrm{size}(\tau_1 \to \tau_2) = 1 + \mathrm{size}(\tau_1) +
\mathrm{size}(\tau_2)$ と定める。

補題を二つ用意する。第一に、$\alpha$ が $\tau$ のどこかに現れるなら、任意の $s$
について $\mathrm{size}(s(\alpha)) \leq \mathrm{size}(s(\tau))$。これは $\tau$ に
関する帰納法で、矢印の場合は出現がどちらの枝にあるかで分ける。第二に、$\alpha$
が矢印型に現れるなら不等号は\emph{狭義}になる。矢印の $\mathrm{size}$ が両枝の和
より $1$ 大きく、どの型の $\mathrm{size}$ も $1$ 以上だからである。

さて $s(\alpha) = s(\tau_1 \to \tau_2)$ を満たす $s$ があったとする。第二の補題は
$\mathrm{size}(s(\alpha)) < \mathrm{size}(s(\tau_1 \to \tau_2))$ を与えるが、
仮定でこの二つは同じ項なので $n < n$ となり矛盾する。
\end{proof}

\begin{corollary}
  $\alpha = \alpha \to \alpha$ は解を持たない。
\end{corollary}

これはちょうど $\lambda x.\,x\,x$ を素朴に単一化したときに現れる方程式である。
occurs check は、そこでアルゴリズムが停止せずに無限に型を作り続けるのを
止めている。

\section{なお残る範囲と限界}
\label{sec:limits}

これは MiniML 実装の検証では\emph{ない}。証明されたのは上の二つの体系に
ついてであって、このリポジトリの SML コードとの隔たりは大きい。次は未着手で
ある。

\begin{itemize}
  \item 単一化アルゴリズムそのもの。occurs check の必要性は第\ref{sec:beyond}節で
  証明したが、単一化の手続きを定義してその健全性と完全性を示してはいない。
  \textsc{I-Lam} は依然として $\tau_1$ を推測し、\textsc{I-App} は関数が既に
  矢印型であることを前提にする。
  \item principal type 定理。
  \item 構文的な一般化と意味論的な一般化の橋渡し。
  \item リスト、組、パターン照合。
  \item SML 実装との対応証明。
\end{itemize}

再帰と多段の安全性は第\ref{sec:beyond}節で埋めた。以前この一覧にあった
「再帰」と「occurs check」は、そこへ移っている。

\subsection{証明が覆っていない反例}

MiniML は比較演算子に
$\forall\alpha.\,\alpha \times \alpha \to \tbool$ という型を与えているが、
関数値は実行時に比較できない。次は型検査を通ってから失敗する。

\begin{lstlisting}
## let f x = x;
f : for all 'a, ('a -> 'a) = <fun>
## f = f;
Runtime error: these values cannot be compared.
\end{lstlisting}

これは現状の MiniML に対する progress の反例である。どちらの Coq ファイルにも
比較演算子は無いので、どちらもこれを排除しない。各ファイルが記述している言語の
範囲では、定理は成立している。MiniML を直すには、Standard ML が
\texttt{''a} で行っているような等価型を入れるか、比較を単相に戻すことになる。

\section{検証方法}

Rocq Prover 9.1.0 で確認した。

\begin{lstlisting}
cd coq/hm-soundness
coqc HMSoundness.v
coqc HMTypeSafety.v
\end{lstlisting}

どちらも終了コード 0 で成功し、\texttt{From Coq} が \texttt{From Stdlib} に
置き換わった旨の deprecation 警告以外に出力は無い。\texttt{Admitted} も
\texttt{admit} も \texttt{Axiom} も無く、38 個の証明はすべて \texttt{Qed} で
閉じている。さらに主要な定理について隠れた仮定の有無を監査した。

\begin{lstlisting}
Print Assumptions hm_infer_sound.
Print Assumptions progress.
Print Assumptions preservation.
\end{lstlisting}

いずれも \texttt{Closed under the global context} と答えるので、追加公理に
依存していない。

本稿は次で組んだ。

\begin{lstlisting}
lualatex -interaction=nonstopmode hm_safety_report_ja.tex
\end{lstlisting}

\end{document}
